\chapter{KẾT LUẬN TẠM THỜI VÀ KẾ HOẠCH TIẾP THEO}
\label{ch:ket-luan}

\emph{Đây là kết luận tính đến thời điểm báo cáo tiến độ. Kết luận định lượng cuối cùng
sẽ được bổ sung sau khi hoàn thành thực nghiệm ở Chương~\ref{ch:ket-qua}.}

\section{Những phần đã hoàn thành}

Tính đến thời điểm này, đồ án đã hoàn thành phần nền tảng của hệ thống:

\begin{itemize}[nosep]
  \item Xác định rõ bài toán, mục tiêu và phạm vi nghiên cứu (Chương~\ref{ch:mo-dau}).
  \item Khảo sát các công trình liên quan và xác định khoảng trống nghiên cứu
        (Chương~\ref{ch:tong-quan}).
  \item Thiết kế kiến trúc hệ thống hai nhánh, quy trình xử lý dữ liệu và giao thức đánh
        giá theo thời gian có kiểm soát rò rỉ (Chương~\ref{ch:phuong-phap}).
  \item Thu thập và kiểm tra dữ liệu: giá của 10 mã VN30 trong giai đoạn 2020 đến quý I
        năm 2026, cùng 10.695 bài tin tức Vietstock có mốc thời gian đến phút và nội dung bài.
  \item Xây dựng và chạy thử quy trình tinh chỉnh PhoBERT cho bài toán phân loại cảm xúc ba lớp.
\end{itemize}

\section{Đối chiếu với mục tiêu nghiên cứu}

Xét theo các mục tiêu đặt ra ở Chương~\ref{ch:mo-dau}, mục tiêu về xây dựng bộ dữ liệu
có thể truy nguyên đã đạt được ở mức chuẩn bị hiện tại, do dữ liệu giá và tin tức đều có
nguồn, mốc thời gian và độ phủ được ghi nhận. Các mục tiêu còn lại, gồm đánh giá PhoBERT
ở tầng cảm xúc, kết hợp hai nguồn dữ liệu, huấn luyện mô hình dự báo và đo phần đóng góp
của cảm xúc, vẫn đang trong giai đoạn thực hiện. Chưa có số liệu để đưa ra kết luận
định lượng cho các mục tiêu này.

Việc chọn PhoBERT làm mô hình ngôn ngữ nền phù hợp với dữ liệu tiếng Việt, tương tự cách
FinBERT được thích nghi cho ngữ liệu tài chính tiếng Anh
\cite{Nguyen2020_200300744v3,Yang2020_200608097v2}. Mức sàn của đồ án là mô hình LSTM hai
nhánh kết hợp giá và cảm xúc, kế thừa ý tưởng từ các nghiên cứu quốc tế nhưng thay đổi
ngôn ngữ, thị trường và mục tiêu phân loại ba lớp
\cite{Halder2022_221107392v1,Gu2024_240716150v1}.

\section{Kế hoạch hoàn thiện}

Các bước còn lại để hoàn thành đồ án gồm:

\begin{enumerate}[nosep]
  \item Tinh chỉnh PhoBERT trên GPU và sinh xác suất cảm xúc cho toàn bộ tập dữ liệu tin tức.
        Dự kiến hoàn thành trong tuần đầu tháng 9.
  \item Tổng hợp cảm xúc theo mã và ngày giao dịch, ghép với đặc trưng giá và tạo nhãn ba lớp.
  \item Chạy các mô hình đối chứng và LSTM hai nhánh, thực hiện thí nghiệm loại bỏ đặc
        trưng có và không có cảm xúc, kèm khoảng tin cậy. Dự kiến hoàn thành trong hai tuần kế tiếp.
  \item Điền kết quả vào Chương~\ref{ch:ket-qua}, hoàn thiện kết luận định lượng và nộp
        trước hạn ngày 23 tháng 9.
\end{enumerate}

Giá trị chính của đồ án không nằm ở giả định rằng cảm xúc chắc chắn cải thiện dự báo, mà
ở việc xây dựng một phép đo có thể phân biệt tín hiệu thực với lợi thế do rò rỉ dữ liệu.
Cách tiếp cận thận trọng này được giữ nhất quán cho tới khi có kết quả cuối cùng.
